% analysis/pipeline.tex

\section{AIM Processing Pipeline}

% ============================================================
\subsection*{Motivation}
% ============================================================

Railway alignment modelling involves multiple domains:
\begin{itemize}
\item global coordinate systems (CRS),
\item geometric alignment representation,
\item physical realization,
\item measurement data,
\item optimization.
\end{itemize}

These domains are often treated separately, leading to fragmented models
and inconsistent results.

Within the AIM framework, these components are unified into a single
processing pipeline.

% ============================================================
\subsection{Operator Chain}
% ============================================================

\begin{definition}[AIM pipeline]
The complete modelling and optimization process is described as a
composition of operators:
\[
\boxed{
\mathcal{P}
=
\mathcal{W}^{-1}
\circ
\mathcal{R}
\circ
\mathcal{M}
\circ
\mathcal{W}
\circ
\mathcal{C}
}
\]
\end{definition}

The operators are:
\begin{itemize}
\item \(\mathcal{C}\): CRS transformation,
\item \(\mathcal{W}\): world-to-track transformation,
\item \(\mathcal{M}\): alignment model (parametric or hybrid),
\item \(\mathcal{R}\): realization operator,
\item \(\mathcal{W}^{-1}\): track-to-world transformation.
\end{itemize}


% ============================================================
\subsection{Detailed Structure}
% ============================================================

The pipeline maps raw data to realized geometry:
\[
\text{CRS data}
\;\xrightarrow{\mathcal{C}}\;
\text{world coordinates}
\;\xrightarrow{\mathcal{W}}\;
(s,d)
\;\xrightarrow{\mathcal{M}}\;
\kappa(s)
\;\xrightarrow{\mathcal{R}}\;
\gamma_{\mathrm{real}}(s)
\;\xrightarrow{\mathcal{W}^{-1}}\;
\text{world geometry}.
\]


% ============================================================
\subsection{Hybrid Model Integration}
% ============================================================

The alignment model operator is hybrid:
\[
\mathcal{M}
:
(\mathbf{T}, \delta\kappa)
\;\mapsto\;
\kappa(s).
\]

Thus, the pipeline operates on both:
\begin{itemize}
\item structured parameters,
\item local corrections.
\end{itemize}


% ============================================================
\subsection{Inverse Problem}
% ============================================================

\begin{definition}[Pipeline inversion]
Given observed data \(x_i\), the inverse problem is:
\[
\min_{\mathbf{x}}
\sum_i
\|
\mathcal{P}(\mathbf{x})(s_i) - x_i
\|^2.
\]
\end{definition}

This corresponds to fitting a model such that the realized geometry
matches observations.


% ============================================================
\subsection{Optimization Embedding}
% ============================================================

The pipeline is embedded into an optimization problem:
\[
\min_{\mathbf{x}}
J\bigl(\mathcal{P}(\mathbf{x})\bigr).
\]

The objective is evaluated on the realized geometry, not on the
parametric model.


% ============================================================
\subsection{Jacobian Structure}
% ============================================================

The derivative of the pipeline follows from the chain rule:
\[
\frac{\partial \mathcal{P}}{\partial \mathbf{x}}
=
D\mathcal{W}^{-1}
\cdot
D\mathcal{R}
\cdot
D\mathcal{M}
\cdot
D\mathcal{W}.
\]

Each component has specific properties:
\begin{itemize}
\item \(D\mathcal{M}\): sparse and structured,
\item \(D\mathcal{R}\): smoothing operator,
\item \(D\mathcal{W}\): nonlinear projection,
\item \(D\mathcal{C}\): linear transformation.
\end{itemize}


% ============================================================
\subsection{Computational Implications}
% ============================================================

The pipeline exhibits:
\begin{itemize}
\item sparsity (model level),
\item nonlinearity (projection level),
\item smoothing (realization level),
\item scaling effects (CRS level).
\end{itemize}

Efficient implementation requires:
\begin{itemize}
\item hybrid representations,
\item block-structured solvers,
\item LOD-dependent evaluation.
\end{itemize}


% ============================================================
\subsection{Level-of-Detail Interpretation}
% ============================================================

Different parts of the pipeline dominate at different scales:
\begin{itemize}
\item large scale: CRS and world geometry,
\item medium scale: parametric alignment,
\item small scale: corrections and realization.
\end{itemize}


% ============================================================
\subsection{Key Insight}
% ============================================================

\begin{proposition}
Alignment modelling is not a single mapping, but a composition of
geometric, physical, and computational operators.
\end{proposition}


% ============================================================
\subsection{Connection to AIM}
% ============================================================

\begin{summary}
The AIM framework is fundamentally defined by its operator pipeline.

It integrates coordinate systems, geometry, physics, and optimization
into a single coherent model.

This unified perspective enables:
\begin{itemize}
\item consistent data integration,
\item physically meaningful optimization,
\item scalable computation.
\end{itemize}

Thus, AIM transforms alignment modelling from a collection of isolated
methods into a structured computational system.
\end{summary}
